\documentclass{article}
\usepackage{fontspec}
\usepackage{polyglossia}
\setmainlanguage{arabic}
\setotherlanguage{english}
\newfontfamily\arabicfont[Script=Arabic]{Amiri}
\setmainfont{Amiri}
\begin{document}
\section*{مشاكل العربية في ملفات PDF}

نسخ النصوص من ملفات PDF العربية ينتج غالباً نصوصاً مشوهة لا يمكن البحث عنها، لأن الملف يحفظ أشكال الحروف وليس الحروف نفسها.

هذا الملف كُتب بأداة قلم (qalam) سنة 2026، وكل كلمة فيه تحمل نصها الأصلي حتى يُنسخ بشكل صحيح: السلام عليكم، لا إله إلا الله، الإسلام.

اللُّغَةُ العَرَبِيَّةُ جَمِيلَةٌ، وَالتَّشْكِيلُ يَجِبُ أَنْ يُنْسَخَ أَيْضاً.

\begin{english}Mixed text works too: the word العربية sits inside an English sentence, and numbers like 3.14 and (12) keep their order.\end{english}

\end{document}
